\begin{figure}[!ht]
\centering
\begin{tikzpicture}[x=1cm,y=1cm,font=\sffamily\small,
  box/.style={draw=reportaccent!65,rounded corners=3pt,
  fill=reportaccent!3,minimum width=3.7cm,minimum height=1.30cm,align=center},
  arr/.style={->,>=stealth,thick,draw=reportaccent}]
\node[box] (input) at (0,0) {Observations\\and goals};
\node[box] (model) at (5.3,0) {Shared model\\and computation};
\node[box] (output) at (10.6,0) {Structured\\decisions};
\draw[arr] (input) -- (model);
\draw[arr] (model) -- (output);
\node[font=\sffamily\bfseries\small,anchor=west] at (-1.85,1.15)
  {An overview of the main idea};
\draw[reportaccent!35] (-1.85,-1.15) -- (12.45,-1.15);
\node[align=center,font=\small,text=black!65] at (5.3,-1.65)
  {Replace this vector placeholder with a full-width method figure or key result.\\
   Use readable labels, consistent colors, and an explanatory caption.};
\end{tikzpicture}
\caption{\textbf{Overview of the proposed approach.} Explain what each component
represents, how information flows through the method, and what readers should
learn from the figure. This is an illustrative placeholder, not a result.}
\label{fig:teaser}
\end{figure}
